\begin{helper}
Let \(s\ge4\), let \(K\) be a finite field of order \(q\), and put
\[
  t=s-2,\qquad
  n=\frac{q^{t+1}-1}{q-1},\qquad
  d_G=\frac{q^t-1}{q-1},
\]
\[
  d_F=n-d_G,\qquad
  \lambda=\sqrt{
    \frac{q^t-1}{q-1}-\frac{q^{t-1}-1}{q-1}}.
\]
Let \(G=G(t,K)\) be the polarity graph of \(\noderef{PolarityGraph}\), and let
\(F\) be the loop-complement of \(G\) in the sense of
\(\noderef{LoopGraphComplement}\).  Then \(G\) is an
\((n,d_G,\lambda)\)-graph, \(F\) is an \((n,d_F,\lambda)\)-graph, and
\((F,G)\) is \(H_s\)-free.
\end{helper}
\begin{proof}
Since \(s\ge4\), the integer \(t=s-2\) satisfies \(t\ge2\), and also
\(s\ge2\).  Applying \(\noderef{PolarityGraphParameters}\) with this \(t\)
and with \(q=|K|\) gives that \(G=G(t,K)\) is an
\((n,d_G,\lambda)\)-graph in the sense of \(\noderef{LoopGraphNdLambda}\),
with exactly the displayed values of \(n\), \(d_G\), and \(\lambda\).

Now \(F\) is the loop-complement of \(G\).  Applying
\(\noderef{LoopGraphComplementNdLambda}\) to the \((n,d_G,\lambda)\)-graph
structure on \(G\) gives that \(F\) is an \((n,n-d_G,\lambda)\)-graph.  By the
definition of \(d_F\), this is precisely that \(F\) is an
\((n,d_F,\lambda)\)-graph.

Finally, since \(s\ge2\), \(\noderef{ComplementPolarityPairHsFree}\) applied
to the same field \(K\) and integer \(s\) says that the pair
\[
  \bigl(\overline{G(s-2,K)},G(s-2,K)\bigr)
\]
is \(H_s\)-free, where the bar denotes loop-complement as in
\(\noderef{LoopGraphComplement}\).  This is exactly the displayed pair
\((F,G)\), because \(t=s-2\).  Hence all three claimed properties hold.
\end{proof}
